\documentclass[12pt, a4paper]{article}

% ── Packages ────────────────────────────────────────────────────────────────
\usepackage[utf8]{inputenc}
\usepackage[T1]{fontenc}
\usepackage{lmodern}
\usepackage[margin=2.5cm]{geometry}
\usepackage{amsmath, amssymb}
\usepackage{graphicx}
\usepackage{hyperref}
\usepackage{natbib}
\usepackage{booktabs}
\usepackage{setspace}
\usepackage{abstract}

\onehalfspacing

% ── Meta-data ────────────────────────────────────────────────────────────────
\title{\textbf{E-CoM Market: An Overview}\\[0.5em]
       \large A Research Paper}

\author{Author Name\\
        \small Institution / Affiliation\\
        \small \texttt{email@example.com}}

\date{\today}

% ── Document ─────────────────────────────────────────────────────────────────
\begin{document}

\maketitle

\begin{abstract}
This paper presents an overview of the E-CoM Market framework.
We describe the motivation, key concepts, and potential applications.
A brief analysis of relevant prior work is included alongside
directions for future research.
\end{abstract}

\tableofcontents
\newpage

% ─────────────────────────────────────────────────────────────────────────────
\section{Introduction}
\label{sec:introduction}

The E-CoM Market represents a novel approach to electronic commerce
modeling. In this paper we outline the core ideas and provide a
structured analysis of the market dynamics involved
\citep{example2023}.

% ─────────────────────────────────────────────────────────────────────────────
\section{Background}
\label{sec:background}

\subsection{Related Work}
Prior research has explored various aspects of electronic commerce
systems \citep{smith2020}. This work builds upon those foundations
while introducing new perspectives on market interactions.

\subsection{Motivation}
The primary motivation for this work stems from the observed gap
in existing models when applied to dynamic, real-time market
environments.

% ─────────────────────────────────────────────────────────────────────────────
\section{Methodology}
\label{sec:methodology}

The methodology follows a structured approach:
\begin{enumerate}
    \item Define the market model formally.
    \item Identify key parameters and variables.
    \item Apply the model to representative case studies.
    \item Evaluate outcomes against established benchmarks.
\end{enumerate}

A summary of the core variables is given in Table~\ref{tab:variables}.

\begin{table}[h!]
    \centering
    \caption{Core model variables}
    \label{tab:variables}
    \begin{tabular}{lll}
        \toprule
        \textbf{Symbol} & \textbf{Description} & \textbf{Units} \\
        \midrule
        $p$  & Market price        & USD \\
        $q$  & Quantity traded     & units \\
        $\lambda$ & Market liquidity & -- \\
        \bottomrule
    \end{tabular}
\end{table}

% ─────────────────────────────────────────────────────────────────────────────
\section{Results}
\label{sec:results}

Preliminary results indicate a strong correlation between market
liquidity and price stability. Further experimental validation
is ongoing.

% ─────────────────────────────────────────────────────────────────────────────
\section{Discussion}
\label{sec:discussion}

The findings suggest that the E-CoM Market model is well-suited
for real-time applications. Limitations include data availability
and the assumptions of the underlying theoretical framework.

% ─────────────────────────────────────────────────────────────────────────────
\section{Conclusion}
\label{sec:conclusion}

This paper introduced the E-CoM Market framework and outlined its
key components and potential applications. Future work will focus
on empirical validation and extension to multi-agent settings.

% ── Bibliography ─────────────────────────────────────────────────────────────
\bibliographystyle{plainnat}
\bibliography{references}

\end{document}
